\chapter{Software Requirement Specification}

The Software Requirement Specification (SRS) defines the operational requirements and constraints of the Intelligent Clinical Document Processing System (ICDPS). Establishing these requirements is necessary to ensure that the implemented system aligns with the objectives of clinical document processing, information extraction, summarization, and SNOMED CT terminology mapping described in previous chapters.

The proposed system integrates multiple processing components within a unified workflow, including document ingestion, OCR, clinical entity extraction, semantic retrieval, and structured output generation. Each component introduces specific requirements related to functionality, performance, reliability, security, and usability. These requirements serve as a reference for subsequent design, implementation, and evaluation activities.

This chapter describes the functional and non-functional requirements of the system together with the hardware, software, and operational conditions considered during development and deployment.

\section{Overall Description}

The Intelligent Clinical Document Processing System (ICDPS) is designed to transform unstructured clinical documents into structured information suitable for review and terminology coding. Users submit clinical documents to the system, which then processes the content through a sequence of extraction, analysis, and coding stages before presenting the resulting information through a review interface.

The system supports the processing of both digitally generated and image-based clinical documents. Its primary outputs include extracted clinical entities, generated summaries, and confidence-ranked SNOMED CT coding suggestions. The architecture is intended to operate as an intermediary processing layer between document sources and downstream healthcare applications.

This section outlines the operating environment, principal system functions, user interactions, and design constraints that define the scope and expected behaviour of the proposed solution.

\subsection{Product Perspective}

The Intelligent Clinical Document Processing System (ICDPS) was designed as a standalone server-side application responsible for processing clinical documents and generating structured outputs. Rather than being developed for a specific Electronic Health Record (EHR) platform, the system operates independently and can be integrated with existing healthcare applications through REST-based interfaces.

Clinical documents are submitted to the processing pipeline, where they undergo document ingestion, text extraction, information extraction, summarization, and SNOMED CT mapping. The resulting entities, summaries, and coding suggestions are then returned to the requesting application through the API. By separating document processing from external healthcare systems, the architecture supports integration with different document repositories, review interfaces, and coding workflows without requiring modifications to the core pipeline.

The system was designed for deployment within a secure healthcare environment. Communication between components is protected using TLS 1.3 encryption, while processed documents and generated outputs are stored in a PostgreSQL database to support auditing, validation, and review activities. All processing is performed within the deployment environment, ensuring that patient information remains within organizational boundaries throughout the workflow.

\subsection{Product Functions}
The ICDPS provides the following functionalities:
\begin{itemize}
    \item \textbf{Document Ingestion:} Accept clinical documents in PDF and image formats through API-based uploads or batch document processing workflows.
    \item \textbf{Text Extraction:} Extract machine-readable text from PDFs using direct extraction; apply OCR to image-based documents.
    \item \textbf{Named Entity Recognition:} Identify and classify clinical entities in extracted text across eight predefined categories.
    \item \textbf{Clinical Summarization:} Generate extractive summaries of clinical documents highlighting key findings.
    \item \textbf{Role-Based Categorization:} Assign extracted entities to clinical role action categories (Pharmacist, Clinician, Radiologist, etc.).
    \item \textbf{SNOMED CT Suggestion:} Retrieve and rank candidate SNOMED CT codes for each extracted entity, presenting the top-5 suggestions with confidence scores.
    \item \textbf{Interactive Review Interface:} Provide a web-based interface for authorized clinicians and coders to review, accept, modify, or reject suggested codes.
    \item \textbf{Audit Logging:} Record all processing events and user actions for compliance and quality assurance purposes.
\end{itemize}

\subsection{User Characteristics}
The ICDPS is intended for two primary user categories:
\begin{itemize}
    \item \textbf{Clinical Coders:} Staff responsible for assigning SNOMED CT codes to clinical documents. Typically possess clinical coding training but may have limited technical expertise. Interact with the system via the web-based review interface.
    \item \textbf{System Administrators:} Technical staff responsible for system configuration, monitoring, and maintenance. Possess technical expertise and access system management functions via command-line and configuration interfaces.
\end{itemize}

\subsection{Constraints}
The following constraints apply to the ICDPS:
\begin{itemize}
    \item \textbf{Data Privacy:} All processing must comply with the UK General Data Protection Regulation (UK GDPR) and NHS Data Security and Protection Toolkit requirements. Personally identifiable patient data must not be transmitted outside the clinical network.
    \item \textbf{Hardware Dependency:} Optimal NER performance requires a CUDA-compatible GPU with at least 8 GB VRAM. CPU-only operation is supported but with reduced inference speed.
    \item \textbf{SNOMED CT Licensing:} SNOMED CT is licensed by SNOMED International. Use of SNOMED CT requires a valid SNOMED CT licence. The system incorporates the UK Edition of SNOMED CT, requiring an NHS licence.
    \item \textbf{Language Limitation:} The current version supports English-language documents only. Multilingual support is identified as a future enhancement.
    \item \textbf{OCR Accuracy:} OCR performance is dependent on document scan quality. Documents with resolution below 150 DPI may exhibit degraded extraction quality.
\end{itemize}

\section{Specific Requirements}

This section outlines the requirements identified for the Intelligent Clinical Document Processing System (ICDPS). These requirements describe the expected functionality of the system, its operational characteristics, and the constraints considered during development. They establish the criteria against which the system is designed, implemented, and evaluated.

\subsection{Functional Requirements}

The functional requirements define the primary capabilities expected from the ICDPS. They describe how the system processes clinical documents, extracts and organizes information, generates SNOMED CT coding suggestions, and presents results for user review. Table~\ref{tab:functional_requirements} summarizes these requirements and their implementation priority.

\begin{longtable}{|p{1.4cm}|p{10.4cm}|p{2cm}|}
\caption{Functional Requirements}
\label{tab:functional_requirements} \\
\hline
\textbf{FR ID} & \textbf{Requirement Description} & \textbf{Priority} \\
\hline
\endfirsthead
\hline
\textbf{FR ID} & \textbf{Requirement Description} & \textbf{Priority} \\
\hline
\endhead
\hline
\endfoot
FR-01 & The system shall accept clinical documents in PDF, JPEG, PNG, and TIFF formats via REST API upload. & High \\
\hline
FR-02 & The system shall extract machine-readable text from PDFs using direct text extraction without OCR. & High \\
\hline
FR-03 & The system shall apply OCR to image-format documents using the Tesseract 5.0 engine to generate UTF-8 text. & High \\
\hline
FR-04 & The system shall automatically detect and classify clinical entities from processed document text into predefined categories such as Diagnosis, Medication, Dosage, Procedure, Anatomy, Allergy, Laboratory Result, and Temporal Expression. & High \\
\hline
FR-05 & The system shall assign a negation status (affirmed/negated/uncertain) to each identified entity. & High \\
\hline
FR-06 & The system shall generate an extractive clinical summary comprising the most clinically significant sentences from the document. & Medium \\
\hline
FR-07 & The system shall assign each extracted entity to a role-specific action category (Pharmacist Action, Clinician Review, Radiologist Review, or General Information). & High \\
\hline
FR-08 & The system shall retrieve the top-5 SNOMED CT concept candidates for each affirmed entity using vector similarity search, ranked by confidence score. & High \\
\hline
FR-09 & The system shall present SNOMED CT suggestions with the concept identifier, preferred term, and confidence score through the review interface. & High \\
\hline
FR-10 & The system shall allow authorized users to accept, modify, or reject SNOMED CT suggestions through the web interface. & High \\
\hline
FR-11 & The system shall log all processing events and user code selection actions with timestamps for audit purposes. & Medium \\
\hline
FR-12 & The system shall support batch processing of multiple documents submitted as a directory archive. & Medium \\
\hline
FR-13 & The system shall return processing results in JSON format through the REST API. & High \\
\hline
FR-14 & The system shall provide an exportable structured report in CSV format summarizing extracted entities and accepted SNOMED CT codes. & Medium \\
\hline
\end{longtable}

\subsection{Non-Functional Requirements}

Non-functional requirements specify the performance and operational characteristics expected from the ICDPS. They address factors such as reliability, security, usability, maintainability, and processing efficiency, which influence the overall quality and dependability of the system. The non-functional requirements considered during system development are presented in Table~\ref{tab:nonfunctional_requirements}.

\begin{longtable}{|p{1.8cm}|p{2.6cm}|p{10cm}|}
\caption{Non-Functional Requirements}
\label{tab:nonfunctional_requirements} \\
\hline
\textbf{NFR ID} & \textbf{Category} & \textbf{Requirement Description} \\
\hline
\endfirsthead
\hline
\textbf{NFR ID} & \textbf{Category} & \textbf{Requirement Description} \\
\hline
\endhead
\hline
\endfoot
NFR-01 & Performance & The system shall complete NER processing of a single-page clinical document within 3 seconds on hardware meeting the minimum specification. \\
\hline
NFR-02 & Throughput & The system shall support batch processing of at least 100 documents per hour during standard operation. \\
\hline
NFR-03 & Accuracy & The NER module shall achieve a token-level F1-score of at least 0.85 on the system test corpus across all eight entity categories. \\
\hline
NFR-04 & SNOMED Precision & The SNOMED CT mapping module shall achieve a precision@1 value of at least 0.85 on the SNOMED test dataset. \\
\hline
NFR-05 & Security & All API communications shall use TLS 1.3 encryption. User authentication shall use JWT tokens with a 24-hour expiry period. \\
\hline
NFR-06 & Availability & The system shall maintain 99\% uptime during scheduled operational hours (08:00--20:00 weekdays). \\
\hline
NFR-07 & Usability & The review interface shall be usable by clinical coders with minimal technical training. \\
\hline
NFR-08 & Maintainability & The system shall support modular architecture and documented APIs for easier updates and maintenance. \\
\hline
NFR-09 & Portability & The system shall be containerized using Docker and deployable on compatible Linux environments. \\
\hline
\end{longtable}

\subsection{Performance Requirements}

Performance requirements define the expected processing speed and responsiveness of the Intelligent Clinical Document Processing System (ICDPS). These requirements establish operational targets for the major processing components and help ensure that clinical documents can be analysed within acceptable time limits.

The NER module is expected to process approximately 500 tokens per second on a system equipped with an NVIDIA A100 GPU. The SNOMED CT candidate retrieval module should return the top five concept suggestions within 200 milliseconds for each entity query. In addition, the web-based review interface should display extraction results for an individual clinical document within 2 seconds on a standard clinical workstation.

The system should maintain stable performance when processing larger documents containing up to 50 pages or approximately 25,000 tokens. These performance targets were selected to support efficient document processing while maintaining a responsive user experience during review and coding activities.

\subsection{Software Requirements}

Software requirements define the software components and technologies used within the ICDPS framework. The selected libraries, frameworks, and development tools support the implementation of document processing workflows, machine learning inference, database operations, and system deployment activities.

Table~\ref{tab:software_requirements} summarizes the software platforms, development environments, and supporting technologies used in the implementation of the Intelligent Clinical Document Processing System (ICDPS).

\begin{longtable}{|p{4.5cm}|p{9cm}|}
\caption{Software Requirements}
\label{tab:software_requirements}\\

\hline
\textbf{Category} &
\textbf{Specification} \\
\hline
\endfirsthead

\hline
\textbf{Category} &
\textbf{Specification} \\
\hline
\endhead

\hline
\multicolumn{2}{r}{Continued on next page} \\
\endfoot

\hline
\endlastfoot


Operating System &
Ubuntu 22.04 LTS (Server), Windows 10+ (Client browser) \\
\hline

Programming Language &
Python 3.10 \\
\hline

NLP Framework &
Hugging Face Transformers 4.38, PyTorch 2.0 \\
\hline

OCR Engine &
Tesseract 5.0 with LSTM models \\
\hline

PDF Processing &
PyMuPDF 1.23 \\
\hline

Vector Search &
FAISS 1.7 \\
\hline

Web Framework &
Flask 3.0 \\
\hline

Database &
PostgreSQL 16 \\
\hline

Containerization &
Docker 24, Docker Compose 2 \\
\hline

Development IDE &
PyCharm / Visual Studio Code \\
\hline

Version Control &
Git / GitHub \\
\hline

\end{longtable}

\subsection{Hardware Requirements}
The performance and efficiency of the proposed system are significantly influenced by the underlying hardware infrastructure used during development and deployment. Hardware requirements define the minimum and recommended computational resources needed for processing clinical documents, executing transformer-based models, and supporting large-scale data operations. Table~\ref{tab:hardware_requirements} presents the hardware specifications required for effective implementation and system execution.

\begin{longtable}{|p{3.6cm}|p{5.5cm}|p{4.2cm}|}
\caption{Hardware Requirements}
\label{tab:hardware_requirements}\\

\hline
\textbf{Component} &
\textbf{Minimum Specification} &
\textbf{Recommended Specification} \\
\hline
\endfirsthead

\hline
\textbf{Component} &
\textbf{Minimum Specification} &
\textbf{Recommended Specification} \\
\hline
\endhead

\hline
\multicolumn{3}{r}{Continued on next page} \\
\endfoot

\hline
\endlastfoot


Processor &
Intel Core i7 (8 cores) &
Intel Xeon (16+ cores) \\
\hline

RAM &
16 GB DDR4 &
64 GB DDR4 \\
\hline

GPU &
NVIDIA GTX 1080 (8 GB VRAM) &
NVIDIA A100 (40 GB VRAM) \\
\hline

Storage &
500 GB SSD &
2 TB NVMe SSD \\
\hline

Network &
1 Gbps Ethernet &
10 Gbps Ethernet \\
\hline

Operating System &
Ubuntu 22.04 LTS &
Ubuntu 22.04 LTS \\
\hline

\end{longtable}

\subsection{Design Constraints}
Design constraints define the technical and operational limitations that influence the architecture, implementation, and deployment of the proposed Intelligent Clinical Document Processing System (ICDPS). The following are the design constraints of ICDPS:
\begin{itemize}
    \item \textbf{Platform Dependency:} The NER module requires CUDA-compatible GPU hardware for production-grade throughput. CPU-only deployment is limited to low-volume testing scenarios.
    \item \textbf{Memory Constraints:} Loading the BioBERT model requires approximately 1.5 GB of GPU VRAM. Simultaneous processing of multiple documents requires proportional VRAM scaling.
    \item \textbf{SNOMED CT Index Size:} The FAISS index for SNOMED CT concept descriptions occupies approximately 3.5 GB of disk space and 6 GB of RAM when loaded.
    \item \textbf{Network Security:} The system must operate within a network that enforces NHS Data Security and Protection Toolkit controls, restricting external API calls.
    \item \textbf{SNOMED CT Licence Compliance:} The system must not distribute SNOMED CT content beyond the licensed deployment site.
\end{itemize}

\section*{Summary}

This chapter outlined the software requirements of the ICDPS, covering functional requirements, non-functional requirements, performance targets, operational constraints, and the supporting technologies required for deployment. These requirements define the expected behaviour and operating environment of the proposed system.

The next chapter focuses on the high-level system architecture and explains how the identified requirements influenced the design of the overall processing framework.
